\label{sec:introduction}

\subsection{The Prison Gerrymandering Problem}

When the Census Bureau conducts the decennial census, it applies a
``usual residence'' rule: each person is counted at the place where they
live and sleep most of the time.\footnote{U.S.\ Census Bureau,
\textit{Residence Criteria and Residence Situations} (2020).}
For most Americans, usual residence is straightforward. For the
approximately 1.4 million people held in state and federal correctional
facilities at the time of the 2020 census, however, the Bureau treats
the prison or jail as the usual residence, counting incarcerated people
at the facility address rather than at their pre-incarceration home.

This counting rule creates what researchers and advocates call
\textit{prison gerrymandering}: when census counts are used to draw
legislative districts, districts containing large correctional facilities
receive an inflated population count, while the communities from which
incarcerated people came receive a deflated count. The practical effect
is representational: rural districts with prisons require fewer actual
non-incarcerated residents to ``fill'' the required equal-population
quota, giving each non-incarcerated resident in those districts more
relative representation than a resident in an urban district without a
prison.

\subsection{Why Prison Gerrymandering Matters for Algorithmic Redistricting}

The \textsc{bisect} algorithmic redistricting pipeline uses Census
population data as its foundational input: population balance is the
primary hard constraint that every district assignment must satisfy.
If the Census data fed into \textsc{bisect} contains the prison
gerrymandering distortion, then the ``equal-population'' districts
\textsc{bisect} produces will be equal only in the nominal Census sense,
not in terms of actual residential population. This has three implications.

First, in states that have enacted prison population adjustment laws,
\textsc{bisect} plans must use the adjusted population data to comply with
state law. Running \textsc{bisect} on unadjusted data in those states
would produce legally non-compliant plans.

Second, the racial and partisan implications of prison gerrymandering
interact directly with the Voting Rights Act analysis. Because incarcerated
populations are disproportionately Black and Hispanic, and because
correctional facilities are concentrated in rural areas, the prison
population distortion systematically reduces the apparent minority
population of urban minority-opportunity districts. This can affect
whether a district clears the majority-minority threshold required for
VRA compliance.

Third, expert witnesses in redistricting proceedings must understand
whether the population data underlying any plan---algorithmic or
enacted---has been adjusted for prison populations, and must be able to
explain the adjustment's implementation and its consequences.

\subsection{Overview}

Section~\ref{sec:background} sets out the constitutional and legal framework
for prison gerrymandering, including the Supreme Court's decision in
\textit{Evenwel v.\ Abbott} (2016) and the survey of state adjustment laws.
Section~\ref{sec:methodology} describes the \textsc{bisect} implementation:
how the \texttt{bisect fetch} command downloads adjusted Census data where
state laws require it, how the adjustment reallocates incarcerated populations
to their home counties using the Census DHC-A file, and what invariants the
adjusted data must satisfy. Section~\ref{sec:results} presents empirical results
for North Carolina, Texas, and California, reporting which districts change
and by how much under prison-adjusted population data. Section~\ref{sec:legal}
addresses the legal implications for expert witnesses and court filings.
Section~\ref{sec:conclusion} concludes.
